\section{Open Questions}
\label{sec:openq}

The following five questions are grounded in what the paper leaves unresolved AND in what
this replication surfaced. They are non-superficial: each is either an untested paper claim
that a small experiment can settle, or a real ambiguity in the paper's genomics-to-clinical
inference chain. Each has concrete next steps.

\subsection*{Q1. Is MBRL-1077 actually \emph{E.\ bugandensis}, or does the 95.26--95.58\%
ANI edge case dissolve under modern species-delineation methods?}
\textbf{Basis.} The paper is emphatic that MBRL-1077 clusters with the other seven strains
as \emph{E.\ bugandensis} (ANI 95.26\%, dDDH 63.9\%). But the modern consensus threshold
for species is ANI $\geq 95\%$ \emph{and} dDDH $\geq 70\%$; MBRL-1077 clears one and fails
the other. Our FastANI reruns land at 95.53--95.58\%. The paper resolves this in favor of
same-species; the SNP tree (Fig.\ 2) contradicts this and shows MBRL-1077 branching
distinctly. This is an unresolved inconsistency \emph{inside the paper}, not just in the
replication. \\
\textbf{Next steps.} Run pyani-blastn ANI + GGDC 3.0 dDDH between MBRL-1077 and (i) the ISS
consensus, (ii) EB-247\textsuperscript{T}, (iii) other \emph{Enterobacter} type strains from
LPSN. Cross-check with core-genome MLST (chewBBACA) and a \texttt{skani} run. If any two of
the four modern methods put MBRL-1077 outside \emph{E.\ bugandensis}, the paper's C8 story
(``MBRL-1077 is a bugandensis carbapenemase producer'') should be reframed as ``a related
species carbapenemase producer'', which changes the epidemiological narrative.

\subsection*{Q2. Are the five ISS isolates truly independent colonization events or one
event that clonally expanded and then was re-sampled five times?}
\textbf{Basis.} Within-ISS ANI is 99.99--100\%; our SNP-count follow-up (not run here)
would test whether the paper's reported ``9, 12, 15, 13, \dots'' filtered-SNP counts are
sufficient to distinguish independent colonization from clonal sampling. Given the isolation
locations (four from WHC, one from ARED foot platform) and a March 2015 window, the
epidemiology matters: five independent colonizations of ISS by \emph{E.\ bugandensis} would
be a stronger microbial-ecology claim than one colonization sampled five times. \\
\textbf{Next steps.} Map Illumina reads of all five ISS isolates against IF3SW-P2 as
reference with bwa-mem2 + bcftools; enumerate strict SNPs (Q30, DP\textgreater 10) outside
repeats and MGEs. Compare with a molecular-clock estimate: at $\sim 10^{-7}$
substitutions/site/generation, does the observed pairwise SNP distance support a common
source $<12$ months prior to sampling, or centuries? If SNPs cluster along a phylogenetic
tree consistent with a single founder + few generations, the ``five independent
colonizations'' framing is wrong.

\subsection*{Q3. Does the ISS \emph{E.\ bugandensis} actually carry a functional MAR
operon --- or does the paper's C7 claim rest on annotation heuristics that a modern
targeted BLAST would revise?}
\textbf{Basis.} Our AMRFinderPlus screen does not call \emph{marA/B/C/R} because it is
scoped to acquired AMR loci. The paper's C7 claim came from RAST subsystem annotation,
which is heuristic and often calls \emph{marRAB} on any \emph{marR}-family regulator hit.
This claim is one of the paper's clinical hooks (``\emph{marA} regulates 60+ genes'') but
was never verified at the level of intact operon structure. \\
\textbf{Next steps.} BLASTn the intact \emph{E.\ coli} K12 MG1655 \emph{marRAB} operon
(nucleotides + $\pm 200$ bp flanks) against all eight assemblies. Confirm (i) all four ORFs
present in the correct order, (ii) intact ribosome binding sites, (iii) intact MarA-box
promoter, (iv) no truncating indels. If any of (i)--(iv) fails, the paper's ``functional
MAR operon'' framing is downgraded to ``a MarR-family regulator homolog'', with different
regulatory implications. Additionally: express \emph{marA} homolog in \emph{E.\ coli}
$\Delta marRAB$ background and assay for the paper's predicted efflux upregulation.

\subsection*{Q4. Does independent 2026-era PathogenFinder actually agree with the paper's
$>79\%$ pathogenicity claim, and is that number meaningful given the reference set has
grown ${\sim}30\times$ since 2018?}
\textbf{Basis.} The paper's entire ``crew health consideration'' framing pivots on
PathogenFinder $>79\%$. We did not re-run PathogenFinder. Between 2018 and 2026 the
PathogenFinder reference (both human and animal pathogens) has grown enormously, and the
tool has been reimplemented (v1 web $\to$ v2 CLI). A modern re-score could easily land
above \emph{or} below 79\%, and either outcome is scientifically interesting: a still-high
score reinforces the paper; a lower score suggests the original number was reference-set
artifact. \\
\textbf{Next steps.} Install PathogenFinder v2 (DTU CGE), score all eight assemblies, and
compare with the paper's five ISS values. Also cross-check with VirulenceFinder + VFDB
BLAST on the same eight assemblies; if the two tools disagree, the paper's confidence
number is not portable and the clinical framing should be rewritten around a range, not a
single point estimate. Report the delta as ``published 2018 79\% $\to$ 2026 X\%'' and treat
that delta as itself a data point about how pathogenicity-prediction ML tools age.

\subsection*{Q5. Do the RefSeq assemblies we (and any 2026+ reader) download match, byte
for byte, the assemblies the paper used --- or has NCBI silently re-scaffolded them?}
\textbf{Basis.} We resolved paper accessions to RefSeq via WGS master ID / strain name;
in one case (EB-247\textsuperscript{T} FYBI $\to$ GCF\_900324475.1) via a strain-name
inference. RefSeq re-scaffolds and re-annotates on its own schedule, and there is no
publicly documented equivalence between the paper's original submissions and the current
GCF accessions. If the sequence changed even at contig-boundary level, the paper's ANI /
dDDH numbers depend on genome pieces that no longer exist under the same names. This is a
general reproducibility concern with all WGS papers, but here it is directly on the
critical path. \\
\textbf{Next steps.} Pull the WGS-original GenBank submissions from ENA/NCBI at their 2018
snapshot (INSDC keeps versioned records), rerun FastANI on \emph{those} against our GCF
downloads, and quantify the sequence-level drift ($\Delta$bp, $\Delta$contigs, $\Delta$N50).
Publish a small manifest (\texttt{assembly\_map.tsv}) mapping ``paper accession'' $\to$
``2018 submission SHA256'' $\to$ ``2026 GCF SHA256''. If drift is non-zero in any of the
eight, add a warning table to REPORT.tex \S\ref{sec:results} and re-render the ANI matrix
on the 2018-frozen assemblies for a true apples-to-apples reproduction of the paper's Table 1.
